\documentclass[12pt, a4paper]{ctexart}
\usepackage[margin=2cm]{geometry}
\usepackage{amsmath, amssymb}
\usepackage{enumitem}
\usepackage{fancyhdr}
\usepackage{xcolor}
\usepackage{tikz}
\usepackage{graphicx}
\usepackage{tabularx}
\usepackage{multicol}

\setCJKmainfont{Hiragino Sans GB}
\setCJKsansfont{Heiti SC}

\pagestyle{fancy}
\fancyhf{}
\fancyhead[L]{培优专题·导数压轴题}
\fancyhead[R]{测试卷}
\fancyfoot[C]{\thepage}
\renewcommand{\headrulewidth}{0.4pt}

\begin{document}

\begin{center}
{\Large\bfseries 培优专题·导数压轴题 测试卷}\\[0.5em]
{\normalsize 满分150分\quad 考试时间120分钟}
\end{center}

\vspace{0.5em}

\noindent\textbf{一、单项选择题（本大题共8小题，每小题5分，共40分。每小题只有一个选项符合题意。）}

\begin{enumerate}[label=\Alph*.]

\item 函数$f(x)=x^3-3x^2+1$在$[-1,4]$上的最大值为（\quad）
\begin{enumerate}[label=]
\item $1$ \qquad \item $-3$ \qquad \item $17$ \qquad \item $-1$
\end{enumerate}

\item 函数$f(x)=\ln x-x$在$(0,+\infty)$上的最大值为（\quad）
\begin{enumerate}[label=]
\item $0$ \qquad \item $-1$ \qquad \item $1$ \qquad \item $1-\ln 2$
\end{enumerate}

\item 函数$f(x)=x^2e^{-x}$（$x\geq 0$）的极大值为（\quad）
\begin{enumerate}[label=]
\item $\dfrac{1}{e}$ \qquad \item $\dfrac{4}{e^2}$ \qquad \item $1$ \qquad \item $\dfrac{2}{e}$
\end{enumerate}

\item 已知$f(x)=ax^3+bx^2+cx+d$，若$f'(0)=0$，$f'(1)=0$，且$f(0)=1$，$f(1)=-1$，则$f(x)$在$(0,1)$内的极值点个数为（\quad）
\begin{enumerate}[label=]
\item $0$ \qquad \item $1$ \qquad \item $2$ \qquad \item 不确定
\end{enumerate}

\item 函数$f(x)=\dfrac{x}{x^2+1}$的值域为（\quad）
\begin{enumerate}[label=]
\item $\left[-\dfrac{1}{2},\dfrac{1}{2}\right]$ \qquad \item $[-1,1]$ \qquad \item $\left(-\dfrac{1}{2},\dfrac{1}{2}\right)$ \qquad \item $(-1,1)$
\end{enumerate}

\item 曲线$y=e^x$在$x=0$处的切线与曲线$y=\ln x$在$x=1$处的切线（\quad）
\begin{enumerate}[label=]
\item 重合 \qquad \item 平行但不重合 \qquad \item 垂直 \qquad \item 相交但不垂直
\end{enumerate}

\item 函数$f(x)=x-\sin x$在$[0,2\pi]$上单调递增的区间为（\quad）
\begin{enumerate}[label=]
\item $[0,\pi]$ \qquad \item $[\pi,2\pi]$ \qquad \item $[0,2\pi]$ \qquad \item $[0,\frac{\pi}{2}]$
\end{enumerate}

\item 函数$f(x)=\ln x-\dfrac{1}{2}x^2+x$的极值点个数为（\quad）
\begin{enumerate}[label=]
\item $0$ \qquad \item $1$ \qquad \item $2$ \qquad \item $3$
\end{enumerate}

\end{enumerate}

\vspace{0.5em}

\noindent\textbf{二、多项选择题（本大题共4小题，每小题5分，共20分。全部选对得5分，部分选对得2分，有选错的得0分。）}

\begin{enumerate}[label=\arabic*.]

\item 关于函数$f(x)=x^3-12x$，下列结论正确的是（\quad）
\begin{enumerate}[label=]
\item $f(x)$在$(-\infty,-2)$上单调递增
\item $f(x)$在$(-2,2)$上单调递减
\item $f(x)$在$(2,+\infty)$上单调递增
\item $f(x)$的图象关于原点对称
\end{enumerate}

\item 已知$f(x)=e^x-1-x$，下列结论正确的是（\quad）
\begin{enumerate}[label=]
\item $f(x)\geq 0$对一切$x\in\mathbb{R}$恒成立
\item $x=0$是$f(x)$的唯一零点
\item 当$x>0$时，$e^x>1+x$
\item 当$x<0$时，$e^x<1+x$
\end{enumerate}

\item 关于函数$g(x)=\ln x+\dfrac{1}{x}$（$x>0$），下列结论正确的是（\quad）
\begin{enumerate}[label=]
\item $g(x)$在$(0,1)$上单调递减
\item $g(x)$在$(1,+\infty)$上单调递增
\item $g(x)$的最小值为$1$
\item $g(x)>0$对$x>0$恒成立
\end{enumerate}

\item 已知函数$h(x)=x\ln x$，下列结论正确的是（\quad）
\begin{enumerate}[label=]
\item $h(x)$在$(0,\frac{1}{e})$上单调递减
\item $h(x)$在$(\frac{1}{e},+\infty)$上单调递增
\item $h(x)$的最小值为$-\dfrac{1}{e}$
\item 当$x>1$时，$h(x)>0$
\end{enumerate}

\end{enumerate}

\vspace{0.5em}

\noindent\textbf{三、填空题（本大题共4小题，每小题5分，共20分。）}

\begin{enumerate}[label=\arabic*.]

\item 曲线$y=x^3-2x+1$在$x=1$处的切线方程为 \underline{\hspace{3cm}}。

\item 函数$f(x)=x^3-3x^2+3x+1$在$[0,2]$上的最大值为 \underline{\hspace{3cm}}。

\item 函数$f(x)=\dfrac{\ln x}{x}$在$[1,e]$上的最大值为 \underline{\hspace{3cm}}。

\item 函数$f(x)=\ln x-ax$（$a>0$）在$[1,e]$上的最大值为$1-a$，则$a$的取值范围为 \underline{\hspace{3cm}}。

\end{enumerate}

\vspace{0.5em}

\noindent\textbf{四、解答题（本大题共6小题，共70分。解答应写出文字说明、证明过程或演算步骤。）}

\begin{enumerate}[label=\arabic*.]

\item （本小题满分10分）\\
已知函数$f(x)=x^3-ax^2+3x$在$x=1$处取得极值。\\
（1）求$a$的值及$f(x)$的单调区间；\\
（2）求$f(x)$在$[-1,3]$上的最大值和最小值。
\vspace{6cm}

\item （本小题满分12分）\\
证明：当$x>0$时，$e^x\geq 1+x+\dfrac{x^2}{2}$。\\
（1）设$g(x)=e^x-1-x-\dfrac{x^2}{2}$，利用导数证明$g(x)\geq 0$；\\
（2）利用(1)的结论证明$e>1+\dfrac{1}{1!}+\dfrac{1}{2!}+\cdots+\dfrac{1}{n!}$对任意正整数$n$成立。
\vspace{6cm}

\item （本小题满分12分）\\
已知函数$f(x)=\ln x-\dfrac{a(x-1)}{x+1}$（$a\in\mathbb{R}$）。\\
（1）当$a=2$时，求$f(x)$的单调区间；\\
（2）当$x>1$时，$f(x)>0$恒成立，求$a$的取值范围。
\vspace{6cm}

\item （本小题满分12分）\\
已知函数$f(x)=e^x-x^2-ax-1$（$a\in\mathbb{R}$）。\\
（1）当$a=0$时，求$f(x)$的最小值；\\
（2）若$f(x)\geq 0$对$x\geq 0$恒成立，求$a$的取值范围；\\
（3）在(2)的条件下，证明：$\displaystyle\sum_{k=0}^{n}\dfrac{1}{k!}\leq e-\dfrac{1}{(n+1)!}$。
\vspace{6cm}

\item （本小题满分12分）\\
已知函数$f(x)=\ln(1+x)-\dfrac{x(1+ax)}{1+x}$（$a\in\mathbb{R}$）。\\
（1）当$a=0$时，求$f(x)$的单调区间；\\
（2）当$a\geq 1$时，证明$f(x)\leq 0$对$x\geq 0$恒成立；\\
（3）利用(2)的结论证明：$\ln\dfrac{n+1}{n}<\dfrac{1}{n}$（$n$为正整数）。
\vspace{6cm}

\item （本小题满分12分）\\
已知函数$f(x)=\ln x$，$g(x)=f(x)-ax$（$a\in\mathbb{R}$）。\\
（1）求$g(x)$的单调区间；\\
（2）证明：$\left(1+\dfrac{1}{n}\right)^n<e<\left(1+\dfrac{1}{n}\right)^{n+1}$（$n$为正整数）；\\
（3）设$a_n=\left(1+\dfrac{1}{n}\right)^n$，证明数列$\{a_n\}$单调递增且有上界。
\vspace{6cm}

\end{enumerate}

\end{document}
